\SourceLeaf{325}%
len, die in den Resultaten nicht mehr erscheinen, hier ohne Übelstand zu anderm Zweck benutzt werden dürfen.

Durch die zwei Punkte der Ellipsoidfläche, auf welche die Aufgabe sich bezieht, werde eine geodaetische Linie, zunächst von unbestimmter Ausdehnung, geführt, und auf derselben ein beliebiger Anfangspunkt gewählt. Das Stück jener Linie von dem Anfangspunkte bis zu einem unbestimmten Punkte werde durch \(u\) bezeichnet; der Winkel, welchen, an letzterm Punkte, die geodaetische Linie mit dem Meridian macht, jene in dem Sinne wachsender \(u\), diesen von Norden nach Süden genommen, durch \(X\); Breite und Länge des unbestimmten Punktes durch \(Y\) und \(Z\). Ich nehme an, dass die Längen von Westen nach Osten, die Azimuthe \(X\) in dem Sinn von Süden nach Westen zu wachsen. Werden nun noch, wie immer bisher, halbe grosse Achse und Excentricität der erzeugenden Ellipse durch \(a\) und \(e\) bezeichnet, so hat man, aus bekannten Gründen
\[
\begin{aligned}
\frac{\mathrm{d}Y}{\mathrm{d}u}&=-\frac{\cos X(1-ee\sin Y^2)^{\frac32}}{(1-ee)a}\quad .\ .\ .\ .\ .\ .\ .\ .\ .\ .\quad(1)\\
\frac{\mathrm{d}Z}{\mathrm{d}u}&=-\frac{\sin X(1-ee\sin Y^2)^{\frac12}}{a\cos Y}\quad .\ .\ .\ .\ .\ .\ .\ .\ .\ .\quad(2)
\end{aligned}
\]
Es ist ferner, nach einem bekannten Lehrsatze die Grösse
\[
\frac{\sin X\cos Y}{\BareRoot{(1-ee\sin Y^2)}}
\]
für alle Punkte derselben geodaetischen Linie constant, und hieraus, wenn man logarithmisch differentiirt
\[
\operatorname{cotang}X\,\mathrm{d}X=(\operatorname{tang}Y-\frac{ee\cos Y\sin Y}{1-ee\sin Y^2})\mathrm{d}Y=\frac{(1-ee)\operatorname{tang}Y}{1-ee\sin Y^2}\mathord{.}\mathrm{d}Y
\]
folglich, aus der Verbindung mit (1),
\[
\frac{\mathrm{d}X}{\mathrm{d}u}=-\frac{\sin X\operatorname{tang}Y(1-ee\sin Y^2)^{\frac12}}{a}\quad .\ .\ .\ .\ .\ .\ .\ .\ .\quad(3)
\]

\indent Wir wollen jedoch unsere Aufgabe allgemeiner fassen, und
\[
\frac{\mathrm{d}X}{\mathrm{d}u}=x,\qquad\frac{\mathrm{d}Y}{\mathrm{d}u}=y,\qquad\frac{\mathrm{d}Z}{\mathrm{d}u}=z
\]
setzen, indem wir zunächst nur voraussetzen, dass \(x,y,z\) irgendwelche gegebene Functionen der beiden Veränderlichen \(X\) und \(Y\) sind. Es entstehe ferner durch neue Differentiation
